\subsection{The basic trigonometric limit}

A central elementary limit is
\[
\lim_{x\to0}\frac{\sin x}{x}=1,
\]
with angles measured in radians.

\begin{interpretationbox}[What the limit says]
Near the origin, the functions \(\sin x\) and \(x\) are approximately equal. Their ratio approaches \(1\). This is the local fact that later produces the derivative of sine.
\end{interpretationbox}

\begin{examplebox}[A simple transformed limit]
Compute
\[
\lim_{x\to0}\frac{\sin(3x)}{x}.
\]
Write
\[
\frac{\sin(3x)}{x}
=3\frac{\sin(3x)}{3x}.
\]
As \(x\to0\), also \(3x\to0\), so
\[
\lim_{x\to0}\frac{\sin(3x)}{x}=3.
\]
\end{examplebox}

\begin{remarkbox}
The purpose is not to build a catalogue of difficult trigonometric limits. One standard pattern is enough to show how a local approximation can become a computational tool.
\end{remarkbox}
